\documentclass[11pt]{article}

\usepackage[a4paper,margin=1in]{geometry}
\usepackage{amsmath,amssymb,amsthm,bm}
\usepackage{booktabs}
\usepackage{array}
\usepackage{tabularx}
\usepackage{longtable}
\usepackage{enumitem}
\usepackage{xcolor}
\usepackage{hyperref}
\usepackage{microtype}

\hypersetup{
    colorlinks=true,
    linkcolor=blue!50!black,
    citecolor=blue!50!black,
    urlcolor=blue!50!black
}

\newtheorem{assumption}{Assumption}
\newtheorem{remark}{Remark}
\newtheorem{definition}{Definition}

\newcommand{\A}{\mathcal A}
\newcommand{\F}{\mathcal F}
\newcommand{\Lcal}{\mathcal L}
\newcommand{\R}{\mathbb R}
\newcommand{\E}{\mathbb E}
\newcommand{\Prob}{\mathbb P}
\newcommand{\norm}[1]{\left\lVert #1 \right\rVert}
\newcommand{\argmax}{\operatorname*{arg\,max}}
\newcommand{\argmin}{\operatorname*{arg\,min}}
\newcommand{\LCB}{\operatorname{LCB}}
\newcommand{\UCB}{\operatorname{UCB}}
\newcommand{\ID}{\textsc{Identify}}
\newcommand{\CERT}{\textsc{Certify}}
\newcommand{\COMMIT}{\textsc{Commit}}

\title{Confidence-Gated Identify-and-Commit Library Bandit\\\large Detailed Algorithmic Specification}
\author{}
\date{}

\begin{document}
\maketitle

\begin{abstract}
This note defines a recurrent-regime bandit algorithm for settings in which the learner repeatedly encounters a finite, initially unknown, collection of operating regimes. The method maintains a library of regime shapes and uses exploration only when a regime must be identified or when the best arm of a regime has not yet been certified. Once the active regime has a certified best arm, the algorithm commits to that arm and monitors its reward for drift. The method is not a UCB or Thompson-sampling policy in stable operation. Confidence intervals are used only as gates: to decide whether a block is stable, whether a regime matches an existing library entry, whether a new regime can be created, and whether one arm is separated enough to justify commitment. The resulting algorithm has three modes, \ID, \CERT, and \COMMIT, and every transition between them is specified explicitly.
\end{abstract}

\section{Setting and notation}

There are $K$ arms,
\begin{equation}
    \A=\{1,\ldots,K\}.
\end{equation}
At time $t$, the learner selects $A_t\in\A$ and observes a reward $Y_t\in\R$. The environment is piecewise stationary. During a stable episode, rewards are generated by an unobserved regime $z$. Regimes may recur: after leaving regime $z$, the system may later return to the same regime.

For a stable regime $z$, the arm means are decomposed as
\begin{equation}
    \mu_k^{(z)}=\ell^{(z)}+\theta_k^{(z)},
    \qquad k\in\A,
    \label{eq:level-shape}
\end{equation}
where $\ell^{(z)}$ is a common reward level and $\bm\theta^{(z)}$ is the action-relevant shape. The shape is normalized by
\begin{equation}
    \sum_{k=1}^K\theta_k^{(z)}=0.
\end{equation}
Only the shape determines the best arm:
\begin{equation}
    k^\star(z)\in\argmax_{k\in\A}\theta_k^{(z)}.
\end{equation}
A common upward or downward shift of all rewards changes $\ell^{(z)}$ but does not change which arm is best.

For any vector $\bm v\in\R^K$, define
\begin{equation}
    P=I-\frac{1}{K}\bm 1\bm 1^\top,
    \qquad
    \phi(\bm v)=P\bm v.
    \label{eq:centering}
\end{equation}
The operator $P$ removes the average coordinate of $\bm v$. The algorithm represents regimes by centered reward shapes, not by absolute reward vectors.

\begin{assumption}[Local sub-Gaussian rewards]
\label{ass:subgaussian}
During a stable episode with active regime $Z_t$, rewards satisfy
\begin{equation}
    Y_t=\mu_{A_t}^{(Z_t)}+\varepsilon_t,
\end{equation}
with
\begin{equation}
    \E[\varepsilon_t\mid \F_{t-1},A_t,Z_t]=0,
\end{equation}
and, for all $\lambda\in\R$,
\begin{equation}
    \E\!\left[\exp(\lambda\varepsilon_t)\mid \F_{t-1},A_t,Z_t\right]
    \leq
    \exp\!\left(\frac{\lambda^2\sigma^2}{2}\right).
    \label{eq:subgaussian}
\end{equation}
\end{assumption}

This assumption is local. It allows different regimes to have different means. The concentration assumption is imposed only on the centered noise around the corresponding regime-arm mean.

\section{Inputs, maintained objects, and modes}

The algorithm has the following inputs:
\begin{itemize}[leftmargin=1.5em]
    \item $K$: number of arms;
    \item $m$: number of samples per arm in one balanced block;
    \item $\sigma$: sub-Gaussian noise proxy;
    \item $\delta\in(0,1)$: confidence level;
    \item $\gamma\geq0$: minimum practically relevant arm gap for commitment;
    \item $\nu,h>0$: tolerance and threshold of the committed-arm drift detector.
\end{itemize}
The parameter $\gamma$ can be set to zero. A positive value prevents commitment when the empirical advantage of the best arm is statistically significant but practically negligible.

The algorithm maintains a library
\begin{equation}
    \Lcal_t=\{1,\ldots,J_t\}.
\end{equation}
Each library entry $j$ stores:
\begin{equation}
    \left(N_j,S_{j,1},\ldots,S_{j,K},\widehat{\bm\theta}^{(j)},a_j\right).
\end{equation}
Here $N_j$ is the number of accepted balanced-block samples per arm assigned to regime $j$, $S_{j,k}$ is the corresponding reward sum for arm $k$, $\widehat{\bm\theta}^{(j)}$ is the estimated centered shape, and $a_j$ is either the certified arm of regime $j$ or the value $\varnothing$ if no arm is certified yet.

The empirical arm means of library entry $j$ are
\begin{equation}
    \widehat\mu_{j,k}=\frac{S_{j,k}}{N_j},
    \qquad k\in\A,
\end{equation}
and its shape estimate is
\begin{equation}
    \widehat{\bm\theta}^{(j)}=P\widehat{\bm\mu}^{(j)},
    \qquad
    \widehat{\bm\mu}^{(j)}=(\widehat\mu_{j,1},\ldots,\widehat\mu_{j,K}).
    \label{eq:library-shape}
\end{equation}
The same number $N_j$ appears for every arm because library entries are updated only with balanced blocks.

The algorithm has exactly three modes:
\begin{equation}
    \ID,
    \qquad
    \CERT,
    \qquad
    \COMMIT.
\end{equation}
In \ID, the algorithm collects balanced blocks until it can associate the current stable shape with an existing regime or start a candidate new regime. In \CERT, the algorithm collects balanced blocks until the best arm of the current regime is certified. In \COMMIT, the algorithm plays one certified arm until the committed-arm detector raises an alarm.

\section{Confidence radii}

Let $n$ denote the number of samples per arm used to estimate a balanced reward vector. Define the coordinate radius
\begin{equation}
    r(n)=\sigma\sqrt{\frac{2\log(8KT/\delta)}{n}}.
    \label{eq:coord-radius}
\end{equation}
A centered shape coordinate is affected both by the error of one coordinate and by the error of the empirical average. Therefore we use the shape radius
\begin{equation}
    b(n)=2r(n).
    \label{eq:shape-radius}
\end{equation}
The exact numerical constant is not important for the algorithmic structure. What matters is that $b(n)$ decreases as more balanced samples are assigned to the same regime.

For a library entry $j$, define the confidence bounds
\begin{equation}
    \LCB_{j,k}=\widehat\theta_{j,k}-b(N_j),
    \qquad
    \UCB_{j,k}=\widehat\theta_{j,k}+b(N_j).
    \label{eq:lcb-ucb}
\end{equation}
These bounds are not used to form a UCB selection index. They are used only to decide whether one arm is sufficiently separated from its competitors.

\section{Balanced blocks and stable-block aggregation}

\begin{definition}[Balanced block]
A balanced block samples every arm exactly $m$ times. Its empirical reward vector is
\begin{equation}
    \widehat v_k
    =
    \frac{1}{m}\sum_{s\in\mathcal B:A_s=k}Y_s,
    \qquad k\in\A,
    \label{eq:block-vector}
\end{equation}
and its centered shape is
\begin{equation}
    \widehat{\bm\psi}=P\widehat{\bm v}.
    \label{eq:block-shape}
\end{equation}
\end{definition}

Balanced blocks are the only observations used to estimate or update full regime shapes. Committed-arm rewards are not used to update $\widehat{\bm\theta}^{(j)}$, because they contain no information about unplayed arms.

The algorithm does not use a fixed periodic probing schedule. Balanced blocks are collected only in \ID\ and \CERT. While in \COMMIT, the algorithm plays only the certified arm.

Stable-block aggregation is defined as follows. The algorithm maintains a temporary candidate aggregate $C$. The aggregate contains $n_C$ samples per arm, reward sums $S^C_k$, empirical means
\begin{equation}
    \widehat\mu^C_k=\frac{S^C_k}{n_C},
\end{equation}
and centered shape
\begin{equation}
    \widehat{\bm\psi}^{C}=P\widehat{\bm\mu}^{C}.
\end{equation}
When a new balanced block $B$ is collected, with shape $\widehat{\bm\psi}^{B}$, it is compatible with the current aggregate $C$ if
\begin{equation}
    \norm{\widehat{\bm\psi}^{B}-\widehat{\bm\psi}^{C}}_\infty
    \leq
    b(m)+b(n_C).
    \label{eq:stable-block}
\end{equation}
If condition~\eqref{eq:stable-block} holds, the block is added to $C$. If it fails, the current aggregate is discarded and a new aggregate is initialized from block $B$.

Thus transition data are not averaged with stable data. A block collected while the environment is moving either fails the compatibility test or starts a new candidate aggregate. This is the only transition-handling mechanism in the algorithm.

\section{Regime identification}

In \ID, the algorithm repeatedly collects balanced blocks and maintains a stable aggregate $C$ as described above. Once $C$ is available, the algorithm compares its shape with the library.

For every stored regime $j\in\Lcal_t$, define
\begin{equation}
    D_j(C)=\norm{\widehat{\bm\psi}^{C}-\widehat{\bm\theta}^{(j)}}_\infty.
    \label{eq:distance}
\end{equation}
The aggregate $C$ is compatible with regime $j$ if
\begin{equation}
    D_j(C)\leq b(n_C)+b(N_j).
    \label{eq:compatible-regime}
\end{equation}
Let
\begin{equation}
    \mathcal M(C)=\left\{j\in\Lcal_t:
    D_j(C)\leq b(n_C)+b(N_j)
    \right\}
    \label{eq:match-set}
\end{equation}
be the match set.

The identification rule is the following.

\paragraph{Unique match.}
If $\mathcal M(C)=\{j\}$, the current regime is identified as $j$. The aggregate $C$ is assigned to library entry $j$, and the library statistics are updated by
\begin{equation}
    S_{j,k}\leftarrow S_{j,k}+S^C_k,
    \qquad
    N_j\leftarrow N_j+n_C,
    \qquad k\in\A.
    \label{eq:library-update}
\end{equation}
The shape $\widehat{\bm\theta}^{(j)}$ is recomputed using~\eqref{eq:library-shape}. The algorithm then checks whether regime $j$ has a certified arm. If yes, it enters \COMMIT. If not, it enters \CERT.

\paragraph{Multiple matches.}
If $|\mathcal M(C)|>1$, the current data do not distinguish the matched regimes. If all regimes in $\mathcal M(C)$ have the same certified arm, the algorithm may commit to that common arm, because the ambiguity is irrelevant for the action. If the matched regimes either have different certified arms or have no certified arm, the algorithm remains in \ID\ and collects more balanced blocks. No library update is performed until the ambiguity is resolved or action-equivalence is established.

\paragraph{No match.}
If $\mathcal M(C)=\varnothing$, the aggregate is treated as a candidate new regime. However, a new library entry is not created immediately. The algorithm enters \CERT\ with candidate $C$. The candidate becomes a new regime only if its best arm is certified and its shape remains separated from all existing library entries. If those conditions are not met, the algorithm keeps collecting balanced blocks and updating the candidate aggregate.

This rule ensures that a regime is never created merely because one block is unmatched. Creation requires both shape separation and best-arm certification.

\section{Best-arm certification}

The purpose of \CERT\ is to decide whether a regime or candidate regime has enough evidence to justify deterministic commitment to one arm.

\subsection{Certification of an existing regime}

For an existing regime $j$, define
\begin{equation}
    \widehat a_j\in\argmax_{k\in\A}\widehat\theta_{j,k}.
\end{equation}
The arm $\widehat a_j$ is certified if
\begin{equation}
    \LCB_{j,\widehat a_j}
    >
    \max_{k\neq \widehat a_j}\UCB_{j,k}+\gamma.
    \label{eq:cert-existing}
\end{equation}
If condition~\eqref{eq:cert-existing} holds, the algorithm sets $a_j=\widehat a_j$ and enters \COMMIT. If $a_j$ was already certified and the condition still holds, no additional certification samples are required; the algorithm commits immediately.

If condition~\eqref{eq:cert-existing} fails, the algorithm collects another balanced block. The block is accepted only if it is stable and compatible with regime $j$. If accepted, it is added to the statistics of $j$, the confidence bounds are recomputed, and the certification test is repeated. If the block is not compatible with regime $j$, certification is aborted and the algorithm returns to \ID\ with a new candidate aggregate started from that block.

Thus certification is not a mandatory phase after every regime match. It is a gate for commitment. A matched regime with an already certified arm can be used immediately.

\subsection{Certification of a candidate new regime}

For a candidate aggregate $C$, define
\begin{equation}
    \widehat a_C\in\argmax_{k\in\A}\widehat\psi^C_k.
\end{equation}
The candidate best arm is certified if
\begin{equation}
    \widehat\psi^C_{\widehat a_C}-b(n_C)
    >
    \max_{k\neq \widehat a_C}\left\{\widehat\psi^C_k+b(n_C)\right\}+\gamma.
    \label{eq:cert-candidate}
\end{equation}
The candidate is separated from the existing library if
\begin{equation}
    D_j(C)> b(n_C)+b(N_j)
    \qquad
    \text{for every }j\in\Lcal_t.
    \label{eq:candidate-separated}
\end{equation}
A new regime is created only when both~\eqref{eq:cert-candidate} and~\eqref{eq:candidate-separated} hold. The new entry $J_t+1$ is initialized by
\begin{equation}
    N_{J_t+1}=n_C,
    \qquad
    S_{J_t+1,k}=S^C_k,
    \qquad k\in\A,
\end{equation}
with shape $\widehat{\bm\theta}^{(J_t+1)}=\widehat{\bm\psi}^{C}$ and certified arm $a_{J_t+1}=\widehat a_C$.

If the candidate arm is not certified, the algorithm collects more balanced blocks. Compatible blocks are added to $C$. Incompatible blocks reset $C$. Therefore a nearly flat candidate, for which $P\widehat{\bm v}\approx\bm 0$, does not become a regime unless one arm is nevertheless certified by~\eqref{eq:cert-candidate}. If all arms are indistinguishable, no committed regime is created from those data.

\section{Commitment and drift detection}

When regime $j$ has a certified arm $a_j$, the algorithm enters \COMMIT. In this mode the arm-selection rule is deterministic:
\begin{equation}
    A_t=a_j.
    \label{eq:commit}
\end{equation}
There is no UCB bonus, no posterior sampling, and no scheduled exploration during \COMMIT.

At the beginning of the committed phase, the algorithm stores a local reward baseline for the committed arm. Let $C_{\mathrm{cur}}$ be the most recent accepted balanced aggregate that led to the current identification or certification. Its local level estimate is
\begin{equation}
    \widehat\ell_{\mathrm{cur}}
    =
    \frac{1}{K}\sum_{k=1}^K\widehat\mu^{C_{\mathrm{cur}}}_k.
\end{equation}
The predicted reward of the committed arm is
\begin{equation}
    \widehat m_{j,a_j}^{\mathrm{cur}}
    =
    \widehat\ell_{\mathrm{cur}}+\widehat\theta_{j,a_j}.
    \label{eq:local-baseline}
\end{equation}
At committed time $t$, define the residual
\begin{equation}
    e_t=Y_t-\widehat m_{j,a_j}^{\mathrm{cur}}.
\end{equation}
The detector statistic is
\begin{equation}
    G_t=\max\{0,G_{t-1}+|e_t|-\nu\},
    \qquad G_{t_0}=0,
    \label{eq:cusum}
\end{equation}
where $t_0$ is the time at which the committed phase starts. A drift alarm is raised when
\begin{equation}
    G_t\geq h.
    \label{eq:alarm}
\end{equation}
After an alarm, the committed phase stops, $G_t$ is reset, and the algorithm returns to \ID. Data collected after the alarm are not assigned to the previous regime unless they are later accepted through the balanced-block identification procedure.

\begin{assumption}[Committed-arm visibility]
\label{ass:visibility}
Whenever the environment leaves a regime $j$ while the algorithm is committed to $a_j$, the mean reward of the committed arm changes by at least
\begin{equation}
    \left|\mu_{a_j}^{(j)}-\mu_{a_j}^{(z')}\right|
    \geq
    \eta_{\min}>0,
    \label{eq:visibility}
\end{equation}
where $z'$ is the new regime.
\end{assumption}

Assumption~\ref{ass:visibility} is the explicit price of eliminating periodic full-shape probing. If a regime change affects only unplayed arms and leaves the committed arm statistically unchanged, the committed-arm detector has no information with which to detect the change.

\section{Complete algorithm}

\begin{table}[t]
\centering
\small
\renewcommand{\arraystretch}{1.25}
\begin{tabularx}{\textwidth}{>{\raggedright\arraybackslash}p{0.17\textwidth} >{\raggedright\arraybackslash}X >{\raggedright\arraybackslash}p{0.18\textwidth}}
\toprule
\textbf{Mode} & \textbf{Action and decision rule} & \textbf{Next mode} \\
\midrule
\ID & Collect balanced blocks and maintain a stable aggregate $C$. Compute $\mathcal M(C)$ from~\eqref{eq:match-set}. If there is a unique match $j$, update regime $j$ using~\eqref{eq:library-update}. If there are multiple action-equivalent matches, use their common certified arm. If there is no match, treat $C$ as a candidate new regime. & \CERT\ or \COMMIT \\
\midrule
\CERT\ for existing regime & Check~\eqref{eq:cert-existing}. If it holds, store $a_j$ and commit. If it fails, collect more balanced blocks compatible with $j$. If a block becomes incompatible, abort certification. & \COMMIT\ if certified; otherwise \CERT\ or \ID \\
\midrule
\CERT\ for candidate regime & Check both candidate certification~\eqref{eq:cert-candidate} and library separation~\eqref{eq:candidate-separated}. Create a new regime only when both hold. Otherwise continue collecting stable compatible blocks. & \COMMIT\ if created; otherwise \CERT \\
\midrule
\COMMIT & Play $A_t=a_j$. Update the committed-arm statistic~\eqref{eq:cusum}. Do not sample other arms. Do not update the full regime shape from committed-arm observations. & Remain in \COMMIT\ unless~\eqref{eq:alarm} holds \\
\midrule
Alarm & If $G_t\geq h$, stop commitment, reset $G_t$, and discard the current post-alarm data for library purposes until new balanced blocks are accepted. & \ID \\
\bottomrule
\end{tabularx}
\caption{Confidence-Gated Identify-and-Commit Library Bandit. Exploration appears only in \ID\ and \CERT. Once an arm is certified, stable operation is deterministic.}
\label{tab:algorithm}
\end{table}

The same procedure can be written as explicit pseudocode.

\begin{center}
\fbox{%
\begin{minipage}{0.94\textwidth}
\textbf{Algorithm: Confidence-Gated Identify-and-Commit Library Bandit}
\begin{enumerate}[leftmargin=1.6em]
    \item Initialize the library $\Lcal=\varnothing$ and set the mode to \ID.
    \item In \ID, collect balanced blocks until a stable aggregate $C$ is available.
    \item Compute the match set $\mathcal M(C)$.
    \item If $\mathcal M(C)=\{j\}$, assign $C$ to regime $j$, update $j$, and check whether $j$ has a certified arm.
    \item If $\mathcal M(C)=\varnothing$, keep $C$ as a candidate new regime and enter \CERT.
    \item If $|\mathcal M(C)|>1$, commit only if all matched regimes have the same certified arm; otherwise continue balanced sampling in \ID.
    \item In \CERT, collect stable balanced blocks until either an existing regime satisfies~\eqref{eq:cert-existing} or a candidate new regime satisfies both~\eqref{eq:cert-candidate} and~\eqref{eq:candidate-separated}.
    \item When a certified arm $a_j$ is available, enter \COMMIT.
    \item In \COMMIT, play $a_j$ at every round and update $G_t$ using~\eqref{eq:cusum}.
    \item If $G_t\geq h$, raise a drift alarm, reset the detector, and return to \ID.
\end{enumerate}
\end{minipage}}
\end{center}

\section{What the algorithm does not do}

The algorithm does not perform periodic probing during committed operation. Therefore it does not guarantee detection of hidden shape changes that are invisible on the committed arm. Such changes are excluded by Assumption~\ref{ass:visibility}.

The algorithm does not create a regime as soon as a block fails to match the library. A new regime is created only after the candidate is stable, separated from all existing regimes, and has a certified best arm.

The algorithm does not require a separate blackout safeguard. If a balanced block has nearly zero centered shape, $P\widehat{\bm v}\approx 0$, then no arm will be certified unless the confidence intervals still separate one arm from the others. In the absence of separation, the algorithm keeps sampling and does not create a usable regime.

\section{Regret interpretation}

The appropriate theoretical statement is a high-probability regret decomposition rather than a stationary-bandit regret bound. Let $G_T$ be the number of true stable episodes up to horizon $T$, let $D_g$ be the delay between the $g$-th true drift and its detection, and let $\Delta_{\max}$ bound one-step regret. On the event that the confidence statements used for matching and certification are valid,
\begin{equation}
    R_T
    \leq
    R_{\mathrm{id}}(T)
    +
    R_{\mathrm{cert}}(T)
    +
    \Delta_{\max}\sum_{g=1}^{G_T}D_g.
    \label{eq:regret}
\end{equation}
A high-probability form adds the failure cost
\begin{equation}
    \delta T\Delta_{\max}.
\end{equation}
The term $R_{\mathrm{id}}(T)$ is the cost of balanced blocks collected to identify regimes. The term $R_{\mathrm{cert}}(T)$ is the cost of balanced blocks collected to certify best arms. The delay term is the cost of continuing to play the previous committed arm until the detector fires.

This decomposition is meaningful even when the number of regime changes grows with $T$. If regime changes occur linearly often, cumulative regret can also grow linearly; the relevant quantity is the cost per identified episode and the delay required to detect a visible drift.

\section{Summary}

The algorithm has a simple interpretation. Balanced blocks identify and certify. Certified regimes are stored in a library. A stored regime can be reused immediately if its best arm is already certified. Stable operation is deterministic: the certified arm is played until its reward becomes inconsistent with the current baseline. No UCB or Thompson-sampling mechanism is used during commitment. Confidence intervals determine when the algorithm is allowed to stop exploring, not which arm should be explored next.

\end{document}
